\documentclass[11pt]{article}

\usepackage{amsmath,amssymb,amsthm}
\usepackage{geometry}
\usepackage{bm}
\usepackage[framemethod=TikZ]{mdframed}

% Table of Contents depth
\setcounter{tocdepth}{2}      % sections and subsections
\setcounter{secnumdepth}{3}   % number down to subsubsections

% Hyperlinks (load last)
\usepackage[
colorlinks=true,
linkcolor=blue,
citecolor=blue,
urlcolor=blue
]{hyperref}
\usepackage{bookmark}

\geometry{margin=1in}

\title{A Geometric Derivation of the Student's $t$ Distribution}
\author{}
\date{}

\begin{document}
	
	\maketitle
	\tableofcontents
	
	\section{Standardizing the sample}
	\begin{mdframed}[
		linewidth=0.8pt,
		roundcorner=6pt,
		backgroundcolor=gray!5,
		linecolor=gray!60,
		innertopmargin=10pt,
		innerbottommargin=10pt,
		innerleftmargin=10pt,
		innerrightmargin=10pt
		]
		
		\subsection*{Step 1. Standardizing the Sample}
		
		Suppose
		
		\[
		X_1,X_2,\ldots,X_n
		\]
		
		are independent random variables satisfying
		
		\[
		X_i\sim N(\mu,\sigma^2),
		\qquad
		i=1,\ldots,n.
		\]
		
		The first step is to remove the unknown location parameter $\mu$ and the
		scale parameter $\sigma$.
		
		For each observation define the standardized random variable
		
		\[
		\boxed{
			Z_i
			=
			\frac{X_i-\mu}{\sigma},
			\qquad
			i=1,\ldots,n.
		}
		\]
		
		Since
		
		\[
		X_i\sim N(\mu,\sigma^2),
		\]
		
		subtracting the mean shifts the distribution without changing its
		variance,
		
		\[
		X_i-\mu
		\sim
		N(0,\sigma^2).
		\]
		
		Next, dividing by the positive constant $\sigma$ rescales the variance.
		
		Indeed,
		
		\[
		E[Z_i]
		=
		E\!\left[
		\frac{X_i-\mu}{\sigma}
		\right]
		=
		\frac{1}{\sigma}
		\left(
		E[X_i]-\mu
		\right)
		=
		0,
		\]
		
		and
		
		\[
		\operatorname{Var}(Z_i)
		=
		\operatorname{Var}
		\left(
		\frac{X_i-\mu}{\sigma}
		\right)
		=
		\frac1{\sigma^2}
		\operatorname{Var}(X_i)
		=
		1.
		\]
		
		Since every affine transformation of a normal random variable is again
		normally distributed, we conclude that
		
		\[
		\boxed{
			Z_i\sim N(0,1),
			\qquad
			i=1,\ldots,n.
		}
		\]
		
		Moreover, because the transformation
		
		\[
		X_i
		\longmapsto
		Z_i
		=
		\frac{X_i-\mu}{\sigma}
		\]
		
		is applied independently to each observation, the independence of the
		sample is preserved.
		
		Therefore,
		
		\[
		Z_1,Z_2,\ldots,Z_n
		\]
		
		remain independent standard normal random variables.
		
		\vspace{0.3cm}
		
		\paragraph{The Standardized Sample Vector}
		
		It is convenient to collect the standardized observations into a single
		column vector,
		
		\[
		\mathbf Z
		=
		\begin{pmatrix}
			Z_1\\
			Z_2\\
			\vdots\\
			Z_n
		\end{pmatrix}.
		\]
		
		Since each coordinate is an independent standard normal random variable,
		the joint density of $\mathbf Z$ is
		
		\[
		\begin{aligned}
			f_{\mathbf Z}(\mathbf z)
			&=
			\prod_{i=1}^{n}
			\frac1{\sqrt{2\pi}}
			\exp\!\left(-\frac{z_i^2}{2}\right)
			\\[0.3cm]
			&=
			(2\pi)^{-n/2}
			\exp\!\left(
			-\frac12
			\sum_{i=1}^{n}z_i^2
			\right).
		\end{aligned}
		\]
		
		Recognizing
		
		\[
		\sum_{i=1}^{n}z_i^2
		=
		\mathbf z^T\mathbf z,
		\]
		
		the density can be written compactly as
		
		\[
		\boxed{
			f_{\mathbf Z}(\mathbf z)
			=
			(2\pi)^{-n/2}
			\exp\!\left(
			-\frac12
			\mathbf z^T\mathbf z
			\right).
		}
		\]
		
		This is precisely the density of the multivariate standard normal
		distribution in $\mathbb R^n$.
		
		Hence,
		
		\[
		\boxed{
			\mathbf Z
			\sim
			N_n(\mathbf0,I_n),
		}
		\]
		
		where
		
		\[
		\mathbf0
		=
		\begin{pmatrix}
			0\\
			\vdots\\
			0
		\end{pmatrix},
		\]
		
		and
		
		\[
		I_n
		=
		\begin{pmatrix}
			1&0&\cdots&0\\
			0&1&\cdots&0\\
			\vdots&\vdots&\ddots&\vdots\\
			0&0&\cdots&1
		\end{pmatrix}
		\]
		
		is the $n\times n$ identity matrix.
		
		This vector formulation is fundamental because it allows us to study the
		sample using the tools of linear algebra. In the next step, we shall apply
		an orthogonal change of coordinates that separates the sample mean from the
		remaining $n-1$ residual degrees of freedom.
		
	\end{mdframed}
	
	\section{Constructing an orthonormal basis}
	\begin{mdframed}[
		linewidth=0.8pt,
		roundcorner=6pt,
		backgroundcolor=gray!5,
		linecolor=gray!60,
		innertopmargin=10pt,
		innerbottommargin=10pt,
		innerleftmargin=10pt,
		innerrightmargin=10pt
		]
		
		\subsection*{Step 2. Constructing an Orthonormal Basis}
		
		Our objective is to decompose the standardized sample vector
		
		\[
		\mathbf Z
		=
		(Z_1,\ldots,Z_n)^T
		\]
		
		into two orthogonal components:
		
		\begin{enumerate}
			\item one representing the sample mean, and
			\item one representing the residual variation around the mean.
		\end{enumerate}
		
		To accomplish this, we construct a new orthonormal coordinate system in
		$\mathbb R^n$.
		
		\paragraph{Step 2.1. The Mean Direction}
		
		Observe that the sample mean is
		
		\[
		\bar Z
		=
		\frac1n
		\sum_{i=1}^{n}Z_i.
		\]
		
		Therefore, the quantity
		
		\[
		\sum_{i=1}^{n}Z_i
		\]
		
		plays a fundamental role.
		
		This suggests introducing the vector
		
		\[
		\mathbf 1
		=
		\begin{pmatrix}
			1\\
			1\\
			\vdots\\
			1
		\end{pmatrix},
		\]
		
		whose dot product with $\mathbf Z$ is precisely
		
		\[
		\mathbf1^T\mathbf Z
		=
		\sum_{i=1}^{n}Z_i.
		\]
		
		The vector $\mathbf1$ is not a unit vector.
		
		Its Euclidean norm is
		
		\[
		\begin{aligned}
			\|\mathbf1\|
			&=
			\sqrt{
				1^2+1^2+\cdots+1^2
			}
			\\
			&=
			\sqrt n.
		\end{aligned}
		\]
		
		Therefore, the corresponding unit vector is
		
		\[
		\boxed{
			e_1
			=
			\frac1{\sqrt n}
			\mathbf1
			=
			\frac1{\sqrt n}
			\begin{pmatrix}
				1\\
				1\\
				\vdots\\
				1
			\end{pmatrix}.
		}
		\]
		
		Indeed,
		
		\[
		\begin{aligned}
			\|e_1\|^2
			&=
			\left(
			\frac1{\sqrt n}
			\right)^2
			\|\mathbf1\|^2
			\\
			&=
			\frac1n(n)
			\\
			&=
			1,
		\end{aligned}
		\]
		
		so that
		
		\[
		\boxed{\|e_1\|=1.}
		\]
		
		Thus, $e_1$ is a unit vector pointing in the direction of equal movement of
		all observations.
		
		\paragraph{Step 2.2. The Orthogonal Complement}
		
		Next, consider all vectors orthogonal to $e_1$.
		
		A vector
		
		\[
		x=(x_1,\ldots,x_n)^T
		\]
		
		is orthogonal to $e_1$ if
		
		\[
		e_1^Tx=0.
		\]
		
		Substituting the definition of $e_1$,
		
		\[
		\frac1{\sqrt n}
		(x_1+x_2+\cdots+x_n)=0.
		\]
		
		Since $\sqrt n>0$, this is equivalent to
		
		\[
		\boxed{
			x_1+x_2+\cdots+x_n=0.
		}
		\]
		
		Thus, every vector orthogonal to $e_1$ has components whose sum is zero.
		
		Geometrically, this subspace represents deviations from the common mean.
		
		Its dimension is
		
		\[
		n-1,
		\]
		
		because one linear constraint has been imposed.
		
		\paragraph{Step 2.3. Completing the Basis}
		
		A fundamental theorem of linear algebra states that every orthonormal set
		can be completed to an orthonormal basis.
		
		Therefore, there exist vectors
		
		\[
		e_2,e_3,\ldots,e_n
		\]
		
		such that
		
		\[
		\boxed{
			\{e_1,e_2,\ldots,e_n\}
		}
		\]
		
		forms an orthonormal basis of $\mathbb R^n$.
		
		Consequently,
		
		\[
		e_i^Te_j=
		\begin{cases}
			1,&i=j,\\
			0,&i\neq j.
		\end{cases}
		\]
		
		Moreover,
		
		\[
		e_2,\ldots,e_n
		\]
		
		all belong to the orthogonal complement of $e_1$, so that
		
		\[
		e_i^Te_1=0,
		\qquad
		i=2,\ldots,n.
		\]
		
		Each of these vectors therefore satisfies
		
		\[
		\sum_{k=1}^{n}(e_i)_k=0.
		\]
		
		They span precisely the residual subspace
		
		\[
		\boxed{
			\left\{
			x\in\mathbb R^n:
			\sum_{i=1}^{n}x_i=0
			\right\}.
		}
		\]
		
		\paragraph{Step 2.4. Forming the Orthogonal Matrix}
		
		Arrange the basis vectors as the rows of a matrix,
		
		\[
		A
		=
		\begin{pmatrix}
			e_1^T\\
			e_2^T\\
			\vdots\\
			e_n^T
		\end{pmatrix}.
		\]
		
		Since the rows are orthonormal,
		
		\[
		AA^T=I_n.
		\]
		
		Likewise,
		
		\[
		A^TA=I_n.
		\]
		
		Hence,
		
		\[
		\boxed{
			A^{-1}=A^T.
		}
		\]
		
		Matrices satisfying this property are called
		\emph{orthogonal matrices}.
		
		Orthogonal matrices preserve inner products, lengths, and angles.
		
		Indeed, for every vector $x\in\mathbb R^n$,
		
		\[
		\begin{aligned}
			\|Ax\|^2
			&=
			(Ax)^T(Ax)
			\\
			&=
			x^TA^TAx
			\\
			&=
			x^TI_nx
			\\
			&=
			x^Tx
			\\
			&=
			\|x\|^2.
		\end{aligned}
		\]
		
		Thus,
		
		\[
		\boxed{
			\|Ax\|=\|x\|.
		}
		\]
		
		This property will be crucial in the next section, where we transform the
		Gaussian vector
		
		\[
		\mathbf Z
		\]
		
		into the new coordinates
		
		\[
		\boxed{
			\mathbf Y=A\mathbf Z.
		}
		\]
		
		The first coordinate of $\mathbf Y$ will represent the sample mean,
		whereas the remaining $n-1$ coordinates will represent the residual
		variation about the mean.
		
	\end{mdframed}
	
	\section{Orthogonal transformations of multivariate normals}
	\begin{mdframed}[
		linewidth=0.8pt,
		roundcorner=6pt,
		backgroundcolor=gray!5,
		linecolor=gray!60,
		innertopmargin=10pt,
		innerbottommargin=10pt,
		innerleftmargin=10pt,
		innerrightmargin=10pt
		]
		
		\subsection*{Step 3. Orthogonal Transformations of Multivariate Normal Random Variables}
		
		In the previous section we constructed an orthogonal matrix
		
		\[
		A
		=
		\begin{pmatrix}
			e_1^T\\
			e_2^T\\
			\vdots\\
			e_n^T
		\end{pmatrix},
		\]
		
		whose rows form an orthonormal basis of $\mathbb R^n$.
		
		We now define a new random vector
		
		\[
		\boxed{
			\mathbf Y
			=
			A\mathbf Z.
		}
		\]
		
		Our objective is to determine the probability distribution of
		$\mathbf Y$.
		
		Since $\mathbf Z$ is multivariate normal,
		
		\[
		\mathbf Z
		\sim
		N_n(\mathbf0,I_n),
		\]
		
		we shall prove that
		
		\[
		\boxed{
			\mathbf Y
			\sim
			N_n(\mathbf0,I_n).
		}
		\]
		
		Consequently, the coordinates
		
		\[
		Y_1,\ldots,Y_n
		\]
		
		remain independent standard normal random variables.
		
		\paragraph{Step 3.1. Recall the Density of $\mathbf Z$}
		
		The density of the standardized sample vector is
		
		\[
		f_{\mathbf Z}(\mathbf z)
		=
		(2\pi)^{-n/2}
		\exp
		\left(
		-\frac12
		\mathbf z^T\mathbf z
		\right).
		\]
		
		Observe that the density depends only on the squared Euclidean norm
		
		\[
		\mathbf z^T\mathbf z
		=
		\|\mathbf z\|^2.
		\]
		
		Thus, the density is rotationally invariant.
		
		\vspace{0.3cm}
		
		\paragraph{Step 3.2. Express the Transformation}
		
		The transformation is
		
		\[
		\mathbf Y=A\mathbf Z.
		\]
		
		Since $A$ is orthogonal,
		
		\[
		A^{-1}=A^T.
		\]
		
		Therefore,
		
		\[
		\boxed{
			\mathbf Z
			=
			A^T\mathbf Y.
		}
		\]
		
		This is a one-to-one linear transformation, so we may apply the
		multivariate change-of-variables formula.
		
		\vspace{0.3cm}
		
		\paragraph{Step 3.3. Compute the Jacobian}
		
		The Jacobian matrix of the inverse transformation is simply
		
		\[
		\frac{\partial\mathbf Z}{\partial\mathbf Y}
		=
		A^T.
		\]
		
		Hence,
		
		\[
		\left|
		\det
		\left(
		\frac{\partial\mathbf Z}
		{\partial\mathbf Y}
		\right)
		\right|
		=
		|\det(A^T)|.
		\]
		
		Using the identity
		
		\[
		\det(A^T)=\det(A),
		\]
		
		we obtain
		
		\[
		=
		|\det(A)|.
		\]
		
		Since $A$ is orthogonal,
		
		\[
		A^TA=I_n.
		\]
		
		Taking determinants,
		
		\[
		\det(A^TA)
		=
		\det(I_n).
		\]
		
		Applying the multiplicative property of determinants,
		
		\[
		\det(A^T)\det(A)=1.
		\]
		
		Since
		
		\[
		\det(A^T)=\det(A),
		\]
		
		we obtain
		
		\[
		(\det A)^2=1.
		\]
		
		Therefore,
		
		\[
		\det(A)=\pm1,
		\]
		
		and consequently,
		
		\[
		\boxed{
			|\det(A)|=1.
		}
		\]
		
		Thus the Jacobian contributes no scaling factor.
		
		\vspace{0.3cm}
		
		\paragraph{Step 3.4. Apply the Change-of-Variables Formula}
		
		The density of $\mathbf Y$ is
		
		\[
		f_{\mathbf Y}(\mathbf y)
		=
		f_{\mathbf Z}(A^T\mathbf y)
		\,
		|\det(A^T)|.
		\]
		
		Since
		
		\[
		|\det(A^T)|=1,
		\]
		
		we obtain
		
		\[
		=
		f_{\mathbf Z}(A^T\mathbf y).
		\]
		
		Substituting the density of $\mathbf Z$,
		
		\[
		=
		(2\pi)^{-n/2}
		\exp
		\left(
		-\frac12
		(A^T\mathbf y)^T
		(A^T\mathbf y)
		\right).
		\]
		
		\vspace{0.3cm}
		
		\paragraph{Step 3.5. Simplify the Quadratic Form}
		
		Using the identity
		
		\[
		(Bx)^T=x^TB^T,
		\]
		
		we have
		
		\[
		(A^T\mathbf y)^T
		=
		\mathbf y^TA.
		\]
		
		Therefore,
		
		\[
		(A^T\mathbf y)^T(A^T\mathbf y)
		=
		\mathbf y^TAA^T\mathbf y.
		\]
		
		Since
		
		\[
		AA^T=I_n,
		\]
		
		it follows that
		
		\[
		=
		\mathbf y^TI_n\mathbf y
		=
		\mathbf y^T\mathbf y.
		\]
		
		Hence,
		
		\[
		\boxed{
			(A^T\mathbf y)^T(A^T\mathbf y)
			=
			\mathbf y^T\mathbf y.
		}
		\]
		
		Geometrically, this identity expresses the fact that orthogonal
		transformations preserve Euclidean length.
		
		\vspace{0.3cm}
		
		\paragraph{Step 3.6. Obtain the Density of $\mathbf Y$}
		
		Substituting the previous identity,
		
		\[
		\begin{aligned}
			f_{\mathbf Y}(\mathbf y)
			&=
			(2\pi)^{-n/2}
			\exp
			\left(
			-\frac12
			\mathbf y^T\mathbf y
			\right).
		\end{aligned}
		\]
		
		This is exactly the density of
		
		\[
		N_n(\mathbf0,I_n).
		\]
		
		Therefore,
		
		\[
		\boxed{
			\mathbf Y
			\sim
			N_n(\mathbf0,I_n).
		}
		\]
		
		\vspace{0.3cm}
		
		\paragraph{Step 3.7. Independence of the Coordinates}
		
		The covariance matrix of $\mathbf Y$ is
		
		\[
		I_n.
		\]
		
		Hence,
		
		\[
		\operatorname{Cov}(Y_i,Y_j)
		=
		\begin{cases}
			1,&i=j,\\
			0,&i\neq j.
		\end{cases}
		\]
		
		A remarkable property of the multivariate normal distribution is that
		uncorrelated components are automatically independent.
		
		Therefore,
		
		\[
		\boxed{
			Y_1,Y_2,\ldots,Y_n
			\text{ are independent }
			N(0,1)
			\text{ random variables.}
		}
		\]
		
		This result is the key step in the proof of Cochran's theorem.
		The orthogonal transformation has separated the original sample into
		independent coordinates while preserving the total squared length,
		
		\[
		\|\mathbf Y\|^2
		=
		\|\mathbf Z\|^2.
		\]
		
		In the next section we shall identify the first coordinate with the
		sample mean and show that the remaining $n-1$ coordinates correspond
		precisely to the residual variation measured by the sample variance.
		
	\end{mdframed}
	
	\section{Interpreting the first coordinate}
	\begin{mdframed}[
		linewidth=0.8pt,
		roundcorner=6pt,
		backgroundcolor=gray!5,
		linecolor=gray!60,
		innertopmargin=10pt,
		innerbottommargin=10pt,
		innerleftmargin=10pt,
		innerrightmargin=10pt
		]
		
		\subsection*{Step 4. Interpreting the First Coordinate}
		
		Having established that
		
		\[
		\mathbf Y=A\mathbf Z,
		\]
		
		where
		
		\[
		\mathbf Z\sim N_n(\mathbf0,I_n),
		\]
		
		and
		
		\[
		A=
		\begin{pmatrix}
			e_1^T\\
			e_2^T\\
			\vdots\\
			e_n^T
		\end{pmatrix},
		\]
		
		our next objective is to determine the statistical meaning of the first
		coordinate of the transformed vector.
		
		Recall that the first basis vector was chosen to be
		
		\[
		e_1
		=
		\frac1{\sqrt n}
		\begin{pmatrix}
			1\\
			1\\
			\vdots\\
			1
		\end{pmatrix}.
		\]
		
		Consequently, the first row of the matrix $A$ is
		
		\[
		e_1^T
		=
		\frac1{\sqrt n}
		\begin{pmatrix}
			1&1&\cdots&1
		\end{pmatrix}.
		\]
		
		\paragraph{Step 4.1. Compute the First Coordinate}
		
		By definition,
		
		\[
		\mathbf Y=A\mathbf Z.
		\]
		
		Writing this matrix product coordinate-wise,
		
		\[
		\mathbf Y
		=
		\begin{pmatrix}
			e_1^T\\
			e_2^T\\
			\vdots\\
			e_n^T
		\end{pmatrix}
		\mathbf Z.
		\]
		
		Therefore,
		
		\[
		Y_1
		=
		e_1^T\mathbf Z.
		\]
		
		Substituting the expression for $e_1$,
		
		\[
		Y_1
		=
		\frac1{\sqrt n}
		\begin{pmatrix}
			1&1&\cdots&1
		\end{pmatrix}
		\begin{pmatrix}
			Z_1\\
			Z_2\\
			\vdots\\
			Z_n
		\end{pmatrix}.
		\]
		
		Carrying out the matrix multiplication,
		
		\[
		\boxed{
			Y_1
			=
			\frac1{\sqrt n}
			\sum_{i=1}^{n}Z_i.
		}
		\]
		
		Thus, the first transformed coordinate is simply the normalized sum of the
		standardized observations.
		
		\vspace{0.3cm}
		
		\paragraph{Step 4.2. Substitute the Definition of $Z_i$}
		
		Recall that each standardized observation is
		
		\[
		Z_i
		=
		\frac{X_i-\mu}{\sigma}.
		\]
		
		Therefore,
		
		\[
		Y_1
		=
		\frac1{\sqrt n}
		\sum_{i=1}^{n}
		\frac{X_i-\mu}{\sigma}.
		\]
		
		Since $\sigma$ is constant,
		
		\[
		=
		\frac1{\sigma\sqrt n}
		\sum_{i=1}^{n}(X_i-\mu).
		\]
		
		Expanding the summation,
		
		\[
		=
		\frac1{\sigma\sqrt n}
		\left(
		\sum_{i=1}^{n}X_i
		-
		\sum_{i=1}^{n}\mu
		\right).
		\]
		
		Because
		
		\[
		\sum_{i=1}^{n}\mu
		=
		n\mu,
		\]
		
		we obtain
		
		\[
		=
		\frac1{\sigma\sqrt n}
		\left(
		\sum_{i=1}^{n}X_i
		-
		n\mu
		\right).
		\]
		
		\vspace{0.3cm}
		
		\paragraph{Step 4.3. Introduce the Sample Mean}
		
		Recall the definition of the sample mean,
		
		\[
		\bar X
		=
		\frac1n
		\sum_{i=1}^{n}X_i.
		\]
		
		Equivalently,
		
		\[
		\sum_{i=1}^{n}X_i
		=
		n\bar X.
		\]
		
		Substituting into the previous expression,
		
		\[
		Y_1
		=
		\frac1{\sigma\sqrt n}
		(n\bar X-n\mu).
		\]
		
		Factor out the common factor $n$,
		
		\[
		=
		\frac{n}{\sigma\sqrt n}
		(\bar X-\mu).
		\]
		
		Since
		
		\[
		\frac{n}{\sqrt n}
		=
		\sqrt n,
		\]
		
		we conclude that
		
		\[
		\boxed{
			Y_1
			=
			\frac{\sqrt n(\bar X-\mu)}
			{\sigma}.
		}
		\]
		
		\vspace{0.3cm}
		
		\paragraph{Step 4.4. Interpretation}
		
		Earlier we proved that every coordinate of the transformed vector
		$\mathbf Y$ is an independent standard normal random variable.
		
		Therefore,
		
		\[
		Y_1
		\sim
		N(0,1).
		\]
		
		Since
		
		\[
		Y_1
		=
		\frac{\sqrt n(\bar X-\mu)}
		{\sigma},
		\]
		
		it follows that
		
		\[
		\boxed{
			\frac{\bar X-\mu}
			{\sigma/\sqrt n}
			\sim
			N(0,1).
		}
		\]
		
		Thus, the first coordinate of the orthogonal transformation is exactly the
		classical standardized sample mean (or $z$-statistic).
		
		Geometrically, the first coordinate measures the projection of the sample
		vector onto the one-dimensional subspace generated by
		
		\[
		(1,1,\ldots,1)^T.
		\]
		
		This direction corresponds to moving all observations equally and therefore
		captures only the common location of the sample.
		
		The remaining $n-1$ coordinates describe movement orthogonal to this
		direction. Since vectors orthogonal to
		
		\[
		(1,1,\ldots,1)^T
		\]
		
		have components whose sum is zero, they represent deviations from the
		sample mean. In the next section, we shall prove that these remaining
		coordinates encode exactly the sample variance.
		
	\end{mdframed}
	
	\section{Interpreting the remaining coordinates}
	\begin{mdframed}[
		linewidth=0.8pt,
		roundcorner=6pt,
		backgroundcolor=gray!5,
		linecolor=gray!60,
		innertopmargin=10pt,
		innerbottommargin=10pt,
		innerleftmargin=10pt,
		innerrightmargin=10pt
		]
		
		\subsection*{Step 5. Interpreting the Remaining Coordinates}
		
		In the previous section we showed that
		
		\[
		Y_1
		=
		\frac{\sqrt n(\bar X-\mu)}{\sigma},
		\]
		
		so the first coordinate of the transformed vector represents the
		standardized sample mean.
		
		Our goal now is to determine the statistical meaning of the remaining
		coordinates
		
		\[
		Y_2,\ldots,Y_n.
		\]
		
		Rather than studying each coordinate individually, we study their squared
		Euclidean length,
		
		\[
		Y_2^2+\cdots+Y_n^2,
		\]
		
		which will turn out to be exactly the standardized residual sum of squares.
		
		\paragraph{Step 5.1. Preservation of Euclidean Length}
		
		Since
		
		\[
		\mathbf Y=A\mathbf Z,
		\]
		
		and the matrix $A$ is orthogonal,
		
		\[
		A^TA=AA^T=I_n,
		\]
		
		orthogonal transformations preserve Euclidean lengths.
		
		Indeed,
		
		\[
		\begin{aligned}
			\|\mathbf Y\|^2
			&=
			(A\mathbf Z)^T(A\mathbf Z)
			\\
			&=
			\mathbf Z^TA^TA\mathbf Z
			\\
			&=
			\mathbf Z^TI_n\mathbf Z
			\\
			&=
			\mathbf Z^T\mathbf Z
			\\
			&=
			\|\mathbf Z\|^2.
		\end{aligned}
		\]
		
		Therefore,
		
		\[
		\boxed{
			\|\mathbf Y\|^2
			=
			\|\mathbf Z\|^2.
		}
		\]
		
		Expanding both norms,
		
		\[
		Y_1^2+Y_2^2+\cdots+Y_n^2
		=
		Z_1^2+Z_2^2+\cdots+Z_n^2.
		\]
		
		Solving for the remaining coordinates,
		
		\[
		\boxed{
			Y_2^2+\cdots+Y_n^2
			=
			Z_1^2+\cdots+Z_n^2-Y_1^2.
		}
		\]
		
		\vspace{0.3cm}
		
		\paragraph{Step 5.2. Express the Right-Hand Side in Terms of the Original Sample}
		
		Recall that
		
		\[
		Z_i
		=
		\frac{X_i-\mu}{\sigma}.
		\]
		
		Hence,
		
		\[
		\sum_{i=1}^{n}Z_i^2
		=
		\frac1{\sigma^2}
		\sum_{i=1}^{n}(X_i-\mu)^2.
		\]
		
		Also,
		
		\[
		Y_1
		=
		\frac{\sqrt n(\bar X-\mu)}{\sigma},
		\]
		
		so
		
		\[
		Y_1^2
		=
		\frac{n(\bar X-\mu)^2}{\sigma^2}.
		\]
		
		Substituting these expressions,
		
		\[
		Y_2^2+\cdots+Y_n^2
		=
		\frac1{\sigma^2}
		\sum_{i=1}^{n}(X_i-\mu)^2
		-
		\frac{n(\bar X-\mu)^2}{\sigma^2}.
		\]
		
		Factoring out the common denominator,
		
		\[
		\boxed{
			Y_2^2+\cdots+Y_n^2
			=
			\frac1{\sigma^2}
			\left[
			\sum_{i=1}^{n}(X_i-\mu)^2
			-
			n(\bar X-\mu)^2
			\right].
		}
		\]
		
		\vspace{0.3cm}
		
		\paragraph{Step 5.3. Apply the Variance Decomposition Identity}
		
		A fundamental algebraic identity states that
		
		\[
		\boxed{
			\sum_{i=1}^{n}(X_i-\mu)^2
			=
			\sum_{i=1}^{n}(X_i-\bar X)^2
			+
			n(\bar X-\mu)^2.
		}
		\]
		
		This identity is often called the
		\emph{decomposition of the total sum of squares}.
		
		Substituting it into the previous equation,
		
		\[
		\begin{aligned}
			Y_2^2+\cdots+Y_n^2
			&=
			\frac1{\sigma^2}
			\left[
			\sum_{i=1}^{n}(X_i-\bar X)^2
			+
			n(\bar X-\mu)^2
			-
			n(\bar X-\mu)^2
			\right]
			\\
			&=
			\frac1{\sigma^2}
			\sum_{i=1}^{n}(X_i-\bar X)^2.
		\end{aligned}
		\]
		
		Thus,
		
		\[
		\boxed{
			Y_2^2+\cdots+Y_n^2
			=
			\frac1{\sigma^2}
			\sum_{i=1}^{n}(X_i-\bar X)^2.
		}
		\]
		
		\vspace{0.3cm}
		
		\paragraph{Step 5.4. Introduce the Sample Variance}
		
		Recall that the sample variance is defined by
		
		\[
		S^2
		=
		\frac1{n-1}
		\sum_{i=1}^{n}(X_i-\bar X)^2.
		\]
		
		Equivalently,
		
		\[
		\sum_{i=1}^{n}(X_i-\bar X)^2
		=
		(n-1)S^2.
		\]
		
		Substituting into the previous equation gives
		
		\[
		\boxed{
			Y_2^2+\cdots+Y_n^2
			=
			\frac{(n-1)S^2}{\sigma^2}.
		}
		\]
		
		\vspace{0.3cm}
		
		\paragraph{Step 5.5. Geometric Interpretation}
		
		The orthogonal transformation separates the standardized sample vector into
		two mutually orthogonal components,
		
		\[
		\mathbf Z
		=
		\underbrace{Y_1e_1}_{\text{mean component}}
		+
		\underbrace{\sum_{i=2}^{n}Y_ie_i}_{\text{residual component}}.
		\]
		
		The first component lies along the direction
		
		\[
		(1,1,\ldots,1)^T,
		\]
		
		which represents a common shift of all observations and therefore contains
		all the information about the sample mean.
		
		The remaining basis vectors
		
		\[
		e_2,\ldots,e_n
		\]
		
		span the orthogonal complement,
		
		\[
		\left\{
		x\in\mathbb R^n:
		\sum_{i=1}^{n}x_i=0
		\right\},
		\]
		
		whose vectors describe deviations from the mean.
		
		The squared length of this residual component is precisely
		
		\[
		\boxed{
			\frac{(n-1)S^2}{\sigma^2}.
		}
		\]
		
		Thus, the sample variance is nothing more than the squared Euclidean length
		of the projection of the sample vector onto the $(n-1)$-dimensional
		subspace orthogonal to the mean direction.
		
	\end{mdframed}
	
	\section{Using orthogonality}
	\begin{mdframed}[
		linewidth=0.8pt,
		roundcorner=6pt,
		backgroundcolor=gray!5,
		linecolor=gray!60,
		innertopmargin=10pt,
		innerbottommargin=10pt,
		innerleftmargin=10pt,
		innerrightmargin=10pt
		]
		
		\subsection*{Step 6. Consequences of Orthogonality}
		
		In the previous sections we established two fundamental facts.
		
		First,
		
		\[
		\mathbf Y
		=
		(Y_1,\ldots,Y_n)^T
		=
		A\mathbf Z,
		\]
		
		where $A$ is an orthogonal matrix.
		
		Second,
		
		\[
		\mathbf Y
		\sim
		N_n(\mathbf0,I_n).
		\]
		
		Since the covariance matrix is the identity,
		
		\[
		\operatorname{Cov}(Y_i,Y_j)
		=
		\begin{cases}
			1,&i=j,\\
			0,&i\neq j.
		\end{cases}
		\]
		
		A special property of the multivariate normal distribution is that
		uncorrelated random variables are automatically independent.
		
		Therefore,
		
		\[
		\boxed{
			Y_1,Y_2,\ldots,Y_n
			\text{ are independent }
			N(0,1)
			\text{ random variables.}
		}
		\]
		
		\vspace{0.3cm}
		
		\paragraph{Step 6.1. The Distribution of the Residual Coordinates}
		
		Recall that we proved
		
		\[
		Y_2^2+\cdots+Y_n^2
		=
		\frac{(n-1)S^2}{\sigma^2}.
		\]
		
		Each random variable
		
		\[
		Y_i,
		\qquad
		i=2,\ldots,n,
		\]
		
		is a standard normal random variable,
		
		\[
		Y_i
		\sim
		N(0,1),
		\]
		
		and these variables are mutually independent.
		
		A fundamental definition of the chi-square distribution states that if
		
		\[
		W_1,\ldots,W_k
		\]
		
		are independent standard normal random variables, then
		
		\[
		W_1^2+\cdots+W_k^2
		\sim
		\chi_k^2.
		\]
		
		Since there are exactly
		
		\[
		n-1
		\]
		
		coordinates
		
		\[
		Y_2,\ldots,Y_n,
		\]
		
		we immediately obtain
		
		\[
		\boxed{
			Y_2^2+\cdots+Y_n^2
			\sim
			\chi_{\,n-1}^2.
		}
		\]
		
		Substituting the identity established in the previous section,
		
		\[
		Y_2^2+\cdots+Y_n^2
		=
		\frac{(n-1)S^2}{\sigma^2},
		\]
		
		gives
		
		\[
		\boxed{
			\frac{(n-1)S^2}{\sigma^2}
			\sim
			\chi_{\,n-1}^2.
		}
		\]
		
		This proves the first statement of Cochran's theorem.
		
		\vspace{0.3cm}
		
		\paragraph{Step 6.2. Independence of the Sample Mean and Sample Variance}
		
		Recall that
		
		\[
		Y_1
		=
		\frac{\sqrt n(\bar X-\mu)}
		{\sigma},
		\]
		
		while
		
		\[
		Y_2^2+\cdots+Y_n^2
		=
		\frac{(n-1)S^2}{\sigma^2}.
		\]
		
		The random variable
		
		\[
		Y_1
		\]
		
		depends only on the first transformed coordinate.
		
		The sample variance depends only on
		
		\[
		Y_2,\ldots,Y_n.
		\]
		
		Since
		
		\[
		Y_1
		\]
		
		is independent of
		
		\[
		(Y_2,\ldots,Y_n),
		\]
		
		it is also independent of every measurable function of those variables.
		
		In particular,
		
		\[
		Y_1
		\]
		
		is independent of
		
		\[
		Y_2^2+\cdots+Y_n^2.
		\]
		
		Therefore,
		
		\[
		\frac{\sqrt n(\bar X-\mu)}{\sigma}
		\]
		
		is independent of
		
		\[
		\frac{(n-1)S^2}{\sigma^2}.
		\]
		
		Finally, multiplying by constants and applying one-to-one transformations
		does not affect independence.
		
		Hence,
		
		\[
		\boxed{
			\bar X
			\quad\text{and}\quad
			S^2
			\text{ are independent.}
		}
		\]
		
		This completes the proof of Cochran's theorem.
		
		\vspace{0.3cm}
		
		\paragraph{Geometric Interpretation}
		
		The orthogonal transformation decomposes the standardized sample vector
		into two orthogonal components,
		
		\[
		\mathbf Z
		=
		Y_1e_1
		+
		\sum_{i=2}^{n}Y_ie_i.
		\]
		
		The first component lies along the one-dimensional subspace generated by
		
		\[
		(1,1,\ldots,1)^T,
		\]
		
		which contains all information about the sample mean.
		
		The remaining $n-1$ coordinates lie in the orthogonal complement,
		
		\[
		\left\{
		x\in\mathbb R^n:
		\sum_{i=1}^{n}x_i=0
		\right\},
		\]
		
		which contains all information about the residual variation.
		
		Because these subspaces are orthogonal, the information about the sample
		mean and the information about the sample variance are completely
		independent. This geometric decomposition is the underlying reason why the
		Student's $t$ statistic exists.
		
	\end{mdframed}
	
	
	\section{Obtaining the chi-square distribution}
	\begin{mdframed}[
		linewidth=0.8pt,
		roundcorner=6pt,
		backgroundcolor=gray!5,
		linecolor=gray!60,
		innertopmargin=10pt,
		innerbottommargin=10pt,
		innerleftmargin=10pt,
		innerrightmargin=10pt
		]
		
		\subsection*{Step 7. Obtaining the Chi-Square Distribution}
		
		Having identified the residual component of the transformed sample vector,
		we now determine its probability distribution.
		
		Recall that
		
		\[
		Y_2,\ldots,Y_n
		\]
		
		are mutually independent random variables satisfying
		
		\[
		Y_i\sim N(0,1),
		\qquad
		i=2,\ldots,n.
		\]
		
		Furthermore, we proved that
		
		\[
		\boxed{
			Y_2^2+\cdots+Y_n^2
			=
			\frac{(n-1)S^2}{\sigma^2}.
		}
		\]
		
		Therefore, the problem reduces to determining the distribution of the sum
		of the squares of independent standard normal random variables.
		
		\paragraph{Step 7.1. Definition of the Chi-Square Distribution}
		
		The chi-square distribution is defined precisely in terms of such sums.
		
		\medskip
		
		\textbf{Definition.}
		If
		
		\[
		W_1,W_2,\ldots,W_k
		\]
		
		are independent standard normal random variables,
		
		\[
		W_i\sim N(0,1),
		\]
		
		then the random variable
		
		\[
		\boxed{
			Q
			=
			W_1^2+W_2^2+\cdots+W_k^2
		}
		\]
		
		is said to follow a chi-square distribution with $k$ degrees of freedom.
		
		We write
		
		\[
		\boxed{
			Q\sim\chi_k^2.
		}
		\]
		
		Thus, the chi-square distribution is simply the distribution of the squared
		Euclidean length of a standard Gaussian vector in $\mathbb R^k$.
		
		\vspace{0.3cm}
		
		\paragraph{Step 7.2. Counting the Degrees of Freedom}
		
		In our setting, the residual component consists of the coordinates
		
		\[
		Y_2,Y_3,\ldots,Y_n.
		\]
		
		The number of these coordinates is
		
		\[
		n-1.
		\]
		
		Since each coordinate is an independent standard normal random variable,
		the definition above applies directly.
		
		Therefore,
		
		\[
		\boxed{
			Y_2^2+Y_3^2+\cdots+Y_n^2
			\sim
			\chi_{\,n-1}^2.
		}
		\]
		
		The appearance of the number
		
		\[
		n-1
		\]
		
		is not accidental.
		
		It reflects the dimension of the residual subspace
		
		\[
		\left\{
		x\in\mathbb R^n:
		\sum_{i=1}^{n}x_i=0
		\right\},
		\]
		
		which has dimension $n-1$.
		
		After estimating the sample mean, one degree of freedom has been used,
		leaving only $n-1$ independent directions in which the observations may
		vary.
		
		\vspace{0.3cm}
		
		\paragraph{Step 7.3. Expressing the Result in Terms of the Sample Variance}
		
		From the previous section,
		
		\[
		Y_2^2+\cdots+Y_n^2
		=
		\frac{(n-1)S^2}{\sigma^2}.
		\]
		
		Since the left-hand side follows a chi-square distribution with
		$n-1$ degrees of freedom, it follows immediately that
		
		\[
		\boxed{
			\frac{(n-1)S^2}{\sigma^2}
			\sim
			\chi_{\,n-1}^2.
		}
		\]
		
		This establishes the first conclusion of Cochran's theorem.
		
		\vspace{0.3cm}
		
		\paragraph{Step 7.4. Geometric Interpretation}
		
		The transformed vector
		
		\[
		\mathbf Y
		=
		(Y_1,\ldots,Y_n)^T
		\]
		
		is obtained from the standardized sample vector by an orthogonal rotation.
		
		The first coordinate,
		
		\[
		Y_1,
		\]
		
		measures the projection of the sample onto the direction
		
		\[
		(1,1,\ldots,1)^T,
		\]
		
		which contains all information about the sample mean.
		
		The remaining coordinates,
		
		\[
		Y_2,\ldots,Y_n,
		\]
		
		span the orthogonal complement of this direction.
		
		Consequently,
		
		\[
		Y_2^2+\cdots+Y_n^2
		\]
		
		is simply the squared Euclidean length of the residual component of the
		sample vector.
		
		Since these residual coordinates form an independent standard Gaussian
		vector in an $(n-1)$-dimensional Euclidean space, their squared length
		follows a chi-square distribution with exactly $n-1$ degrees of freedom.
		
		Thus, the degrees of freedom of the chi-square distribution arise
		naturally from the dimension of the residual subspace.
		
	\end{mdframed}
	
	\section{Proving independence}
	\begin{mdframed}[
		linewidth=0.8pt,
		roundcorner=6pt,
		backgroundcolor=gray!5,
		linecolor=gray!60,
		innertopmargin=10pt,
		innerbottommargin=10pt,
		innerleftmargin=10pt,
		innerrightmargin=10pt
		]
		
		\subsection*{Step 8. Proving the Independence of the Sample Mean and Sample Variance}
		
		The final statement of Cochran's theorem asserts that the sample mean and
		the sample variance are independent random variables.
		
		This result is highly remarkable because it is a special property of the
		normal distribution. For most probability distributions, the sample mean
		and sample variance are \emph{not} independent.
		
		We now prove this result using the orthogonal decomposition developed in
		the previous sections.
		
		\paragraph{Step 8.1. Recall the Orthogonal Transformation}
		
		We introduced the orthogonal transformation
		
		\[
		\mathbf Y
		=
		A\mathbf Z,
		\]
		
		where
		
		\[
		\mathbf Z
		\sim
		N_n(\mathbf0,I_n).
		\]
		
		Since orthogonal transformations preserve the multivariate standard normal
		distribution,
		
		\[
		\mathbf Y
		\sim
		N_n(\mathbf0,I_n).
		\]
		
		Therefore,
		
		\[
		\boxed{
			Y_1,Y_2,\ldots,Y_n
			\text{ are independent }
			N(0,1)
			\text{ random variables.}
		}
		\]
		
		Thus, every coordinate of the transformed vector is independent of all the
		others.
		
		\paragraph{Step 8.2. Express the Sample Mean in Terms of $Y_1$}
		
		Earlier we proved that
		
		\[
		\boxed{
			Y_1
			=
			\frac{\sqrt n(\bar X-\mu)}
			{\sigma}.
		}
		\]
		
		Solving this equation for the sample mean,
		
		\[
		\bar X-\mu
		=
		\frac{\sigma}{\sqrt n}Y_1,
		\]
		
		and therefore,
		
		\[
		\boxed{
			\bar X
			=
			\mu
			+
			\frac{\sigma}{\sqrt n}Y_1.
		}
		\]
		
		Hence, the sample mean depends only on the single random variable
		\(Y_1\).
		
		\paragraph{Step 8.3. Express the Sample Variance in Terms of the Remaining Coordinates}
		
		We also proved that
		
		\[
		Y_2^2+\cdots+Y_n^2
		=
		\frac{(n-1)S^2}{\sigma^2}.
		\]
		
		Solving for \(S^2\),
		
		\[
		\boxed{
			S^2
			=
			\frac{\sigma^2}{n-1}
			\left(
			Y_2^2+\cdots+Y_n^2
			\right).
		}
		\]
		
		Thus, the sample variance depends only on the coordinates
		
		\[
		Y_2,\ldots,Y_n.
		\]
		
		Notice that the coordinate \(Y_1\) does not appear in this expression.
		
		\paragraph{Step 8.4. Independence is Preserved Under Functions}
		
		Since
		
		\[
		Y_1
		\]
		
		is independent of
		
		\[
		(Y_2,\ldots,Y_n),
		\]
		
		a fundamental theorem of probability states that any measurable function of
		\(Y_1\) is independent of any measurable function of
		
		\[
		(Y_2,\ldots,Y_n).
		\]
		
		More precisely, if
		
		\[
		U=g(Y_1),
		\]
		
		and
		
		\[
		V=h(Y_2,\ldots,Y_n),
		\]
		
		for measurable functions \(g\) and \(h\), then
		
		\[
		U
		\quad\text{and}\quad
		V
		\]
		
		are independent.
		
		In our case,
		
		\[
		\bar X
		=
		\mu
		+
		\frac{\sigma}{\sqrt n}Y_1
		=
		g(Y_1),
		\]
		
		while
		
		\[
		S^2
		=
		\frac{\sigma^2}{n-1}
		\left(
		Y_2^2+\cdots+Y_n^2
		\right)
		=
		h(Y_2,\ldots,Y_n).
		\]
		
		Therefore,
		
		\[
		\boxed{
			\bar X
			\text{ and }
			S^2
			\text{ are independent.}
		}
		\]
		
		\paragraph{Step 8.5. Conclusion of Cochran's Theorem}
		
		Combining the results obtained throughout the proof, we have established
		that
		
		\[
		\boxed{
			\frac{\bar X-\mu}
			{\sigma/\sqrt n}
			\sim
			N(0,1),
		}
		\]
		
		\[
		\boxed{
			\frac{(n-1)S^2}{\sigma^2}
			\sim
			\chi^2_{\,n-1},
		}
		\]
		
		and
		
		\[
		\boxed{
			\bar X
			\quad\text{and}\quad
			S^2
			\text{ are independent.}
		}
		\]
		
		These three results together constitute the \emph{Normal Sampling Theorem},
		also known as \emph{Cochran's Theorem}. They provide the theoretical
		foundation for the construction of Student's \(t\)-statistic,
		
		\[
		T
		=
		\frac{\bar X-\mu}
		{S/\sqrt n},
		\]
		
		which can be written as
		
		\[
		T
		=
		\frac{
			\displaystyle
			\frac{\bar X-\mu}
			{\sigma/\sqrt n}
		}{
			\displaystyle
			\sqrt{
				\frac{(n-1)S^2/\sigma^2}
				{n-1}
			}
		}.
		\]
		
		Since the numerator is a standard normal random variable, the denominator
		is the square root of an independent chi-square random variable divided by
		its degrees of freedom, it follows that
		
		\[
		\boxed{
			T
			\sim
			t_{\,n-1}.
		}
		\]
		
		Thus, the Student's \(t\)-distribution emerges naturally from the
		orthogonal decomposition of a multivariate normal sample into two
		independent components: one describing the sample mean and the other
		describing the residual variation about that mean.
		
	\end{mdframed}
	
\end{document}